\section{Unconditional reductions}\label{sec:reductions}

This section contains the paper's unconditional results; neither uses a coding hypothesis. Both reduce adaptive
machines to oblivious ones. \Cref{thm:LA} shows that the uniform exclusions of the ninth paper~\cite{SourceCovers2026},
which are proved there for one fixed computation geometry, apply to every machine of sufficiently low adaptivity: such a
machine collapses to the oblivious case. \Cref{thm:EM} works at block scale. There every linear-time transposer,
adaptive or not, yields one external-memory transposition with linearly many I/Os that is oblivious on a set of inputs of
almost full entropy. Consequently an I/O lower bound of order $\omega(d)$ for this oblivious form of transposition, proved
without the indivisibility assumption of the classical external-memory bounds, would decide~(T) negatively
(\cref{cor:EM}).

\paragraph{Machines and traces.} We use the machine model of~\cite[Section~2.1]{BoundedNesting2026}. A machine~$M$ is a
deterministic Turing machine with a finite state set~$Q$, a read-only two-way input tape, $k\ge1$ two-way work tapes over
a finite alphabet~$\Gamma$ that contains the blank, and a write-only one-way output tape that receives at most one bit per
step. All heads start at cell~$0$, every head moves by $-1$, $0$ or $+1$ per step, every step writes the cell under each
work head, and the input cells outside $[0,L)$ hold an end marker. The \emph{trace} $\tau(X)$ of the run on an input~$X$
is the sequence of all head moves together with the append times and the halting time; for $X$ uniform on $\{0,1\}^L$,
the \emph{probability} of a trace $\tau_0$ is $\Pr[\tau(X)=\tau_0]$. The \emph{value} $v_t(X)$ of step~$t$ consists of
the state after the step, the $k$ symbols written and the output bit appended at the step (or a symbol for ``none''). It
takes at most $2^{\beta(M)}$ values, where $\beta(M)=\lceil\log_2\card Q\rceil+k\lceil\log_2\card\Gamma\rceil+2$ is the
\emph{value width} of~$M$~\cite[Section~4.1]{FourthPaper2026}. We say that $M$ \emph{transposes at dimension~$d$ within
$c(L+1)$ steps} if on every input $X\in\{0,1\}^L$, $L=d^2b(d)$, it halts within $c(L+1)$ steps with output~$X^\T$.

\subsection{Low adaptivity}\label{sec:reductions-LA}

We recall the return families of the ninth paper. Let $x_1,\dots,x_k$ be the work-head walks of a trace on the times
$0,\dots,T-1$. A \emph{long hop} of head~$h$ is a pair $(p,u)$ of times with $u-p\ge2$, $x_h(p)=x_h(u)$ and
$x_h(t)\ne x_h(u)$ for $p<t<u$~\cite[Section~2.2]{BoundedNesting2026}. Its \emph{side} is upper or lower according as the
head stays above or below the cell $x_h(u)$ between $p$ and~$u$, and its \emph{owner} is the first time in $(p,u)$ at which
the head is farthest from that cell. A \emph{complete family} is the set of all long hops of one head, on one side, with
one owner; its \emph{hop count} is its number of hops, and its \emph{full span} is the interval $(p_G,U_G]$ of times, where
$(p_G,U_G)$ is its outermost hop~\cite[Definition~2.1]{SourceCovers2026}. The complete families of the trace admit an
\emph{$N$-class span partition} if they can be partitioned into $N$ classes whose full spans are pairwise disjoint within
each class, and a set of complete families is a \emph{$\mathrm Q_2$ core} if it admits such a partition into two
classes~\cite[Definition~2.1]{SourceCovers2026}. An $N$-class span partition exists exactly when no time lies in more than
$N$ full spans: one direction is clear, and for the other, colour the families greedily in the order of the left ends of
their spans; when a family is coloured, the families already coloured whose spans meet its span all contain the first time
of its span, so they use at most $N-1$ colours.

\begin{theorem}[Theorem LA: low adaptivity]\label{thm:LA}
\notformalized
Let $M$ be a machine with value width $\beta=\beta(M)$, let $c>0$, and let $N\ge1$ and $q\ge0$ be integers.
\begin{enumerate}[label=(\alph*)]
\item Let $\eps>0$ and $d_0$ be the constants that~\cite[Corollary~4.5]{SourceCovers2026} provides for the parameters
$(2c,\beta,N,q)$, and let $0\le\eta<\eps$. Then at no $d\ge d_0$ does $M$ have both of the following properties at
$L=d^2b(d)$:
\begin{enumerate}[label=(\roman*)]
\item $M$ transposes at dimension~$d$ within $c(L+1)$ steps;
\item some trace $\tau_0$ has probability at least $2^{-\eta L}$, its complete families admit an $N$-class span partition,
and at most $q$ of them have more than $\eps L$ hops.
\end{enumerate}
\item Let $\eps>0$ and $d_0$ be the constants that~\cite[Corollary~4.3]{SourceCovers2026} provides for the parameters
$(2c,\beta,N)$, and let $0\le\eta<\eps$. Then the conclusion of~(a) holds when the last clause of~(ii) is replaced by:
the complete families of $\tau_0$ split into a $\mathrm Q_2$ core, at most one further family, and families of at most
$\eps L$ hops each.
\end{enumerate}
\end{theorem}

\begin{informalproof}
We prove~(a). Suppose that $M$ has properties (i) and~(ii) at some $d\ge d_0$, and let $W$ be the set of inputs
$X\in\{0,1\}^L$ with $\tau(X)=\tau_0$. Then
\begin{equation}\label{eq:LA-fibre}
 \card W=2^L\Pr[\tau(X)=\tau_0]\ge2^{(1-\eta)L}.
\end{equation}
We show that $\tau_0$ and~$W$ form an admissible geometry in the sense of~\cite[Definition~2.2]{SourceCovers2026}. This
is the modelling step by which the ninth paper applies its theorems to computations; its fixed geometry ``can be an
oblivious schedule or one trace shared by a set of inputs''~\cite[Section~2.1]{SourceCovers2026}.

\emph{The geometry.} The trace $\tau_0$ fixes the halting time~$T$, the input-head walk~$x_0$ and the work-head walks
$x_1,\dots,x_k$ on $\{0,\dots,T-1\}$, all of which move by at most one cell per step, and the append times. By~(i), the
machine appends exactly $L$ bits on every $X\in W$; since it appends at most one bit per step, the append times are
$t_0<t_1<\dots<t_{L-1}<T$, and output position~$r$ is appended at step~$t_r$. Moreover $T\le c(L+1)\le2cL$, because
$L\ge1$. Let $x(t)$ be $x_0(t)$ clamped to $[0,L-1]$.

\emph{Local determination.} Let $X\in W$ and $t<T$. Step~$t$ applies the transition function of~$M$ to the state at
time~$t$, which is the initial state if $t=0$ and the state component of $v_{t-1}(X)$ otherwise; to the input symbol, which
is $X_{x_0(t)}$ if $0\le x_0(t)<L$ and the end marker otherwise; and to the symbols under the work heads. All runs on~$W$
move their heads as $\tau_0$ prescribes, and every step writes the cell under each work head. Hence the cell under work
head~$h$ at time~$t$ holds the symbol written at step $\prev_h(t)=\max\{u<t:x_h(u)=x_h(t)\}$, a component of
$v_{\prev_h(t)}(X)$, and the blank if that set is empty. This is the local recursion
of~\cite[Lemma~4.2]{FourthPaper2026}, which has a written proof there; its proof uses only the head positions, so it holds
for all inputs that follow~$\tau_0$, as in the proof of~\cite[Corollary~6.4]{BoundedNesting2026}. If
$\prev_h(t)\ne t-1$, then $(\prev_h(t),t)$ is a long hop of head~$h$. Consequently, for $X,Y\in W$: if
$X_{x(t)}=Y_{x(t)}$ and $v_u(X)=v_u(Y)$ for $u=t-1$ (when $t\ge1$) and for the lower ends~$u$ of all long hops ending
at~$t$, then $v_t(X)=v_t(Y)$. These parents are the state edge and the work previous-visit edges
of~\cite[Section~2.1]{SourceCovers2026}, and the input bit at the clamped position $x(t)$ is the direct read there; when
$x_0(t)$ lies outside the input, that read is superfluous, which only enlarges the source sets, as noted there.

\emph{Correctness and capacity.} For $X\in W$ and $0\le u<L$, the output bit at position $\tpos(u)$ is appended at step
$t_{\tpos(u)}$ and is a component of $v_{t_{\tpos(u)}}(X)$; by~(i) it equals~$X_u$. This is the transpose demand
of~\cite[Section~2.1]{SourceCovers2026}. The records $v_t$ take at most $2^\beta$ values.

Thus the walks and append times of~$\tau_0$, the correctness set~$W$, the values~$v_t$ as records and the map that
extracts the output bit from a record form an admissible geometry with entry width $b(d)$, runtime constant~$2c$, capacity
$\beta$ and the $N$-class span partition of~(ii). Its core and exceptional families are labels only; for~(a) we take both
empty and let $a_0$ be the largest hop count of any family, or $0$ if there is none. Its correctness deficit is
$\Delta=L-\log_2\card W\le\eta L$, by~\eqref{eq:LA-fibre}.

Let $\mathcal E$ be the set of complete families of~$\tau_0$ with more than $\eps L$ hops. By~(ii), $\card{\mathcal E}\le q$.
Every other family has at most $\eps L$ hops, hence at most $a=\lfloor\eps L\rfloor$, and $a\le\eps L$. Since $d\ge d_0$,
\cite[Corollary~4.5]{SourceCovers2026}, which is Lean-verified there and permits a skeleton~$\mathcal E$ that differs from
the geometry's labels, gives $\Delta\ge\eps L$. Hence $\eps L\le\eta L$, contradicting $\eta<\eps$.

For~(b), label the given $\mathrm Q_2$ core as the core, the at most one further family as the exceptional set, and put
$a_0=\lfloor\eps L\rfloor$, which bounds the hop count of every remaining family. \Cref{eq:LA-fibre} and the same
modelling give an admissible geometry with one exception and $a_0\le\eps L$, and \cite[Corollary~4.3]{SourceCovers2026},
which is Lean-verified there, gives $\Delta\ge\eps L$. This contradicts $\Delta\le\eta L$ as before.
\end{informalproof}

Hypothesis~(ii) holds, for instance, if at most $2^{\eta L}$ distinct traces occur and each of them has the stated family
structure. More generally it holds if the trace has Shannon entropy $H(\tau(X))\le\eta L$ for a uniformly random input and
every trace of probability at least $2^{-\eta L}$ has that structure, since
$H(\tau(X))\ge-\log_2\max_{\tau}\Pr[\tau(X)=\tau]$.

The theorem is the freezing argument of the proof of~\cite[Corollary~6.4]{BoundedNesting2026} applied to the unconditional
exclusions of the ninth paper: the information of the trace is paid once, through the probability of~$\tau_0$, and it is
smaller than the deficit $\eps L$ that the frozen geometry forces. It thereby moves the classes excluded by the ninth paper
from oblivious machines to machines of low adaptivity, without a coding hypothesis.

\begin{remark}[Scope of Theorem~LA]\label{rem:LA-scope}
\begin{enumerate}[label=(\arabic*)]
\item The constant~$\eps$ is not explicit. The ninth paper obtains its uniform constants by contradiction, extracting
convergent hop profiles by compactness~\cite[Section~4.1]{SourceCovers2026}. \Cref{thm:LA} therefore excludes adaptivity
below an unspecified threshold, whereas the conditional Corollary~F$''$ of the eighth
paper~\cite[Corollary~6.4]{BoundedNesting2026} allows every $\eta<\frac12$.
\item The modelling step in the proof is a written argument. The formal object of the ninth paper packages the
fixed-geometry data as hypotheses and does not derive them from a Turing machine~\cite[Section~2.1]{SourceCovers2026};
the identification above is the one by which that paper applies its theorems to oblivious machines, used here for one
frozen trace.
\item The proof uses~\cite[Corollary~4.5]{SourceCovers2026} only through its bound on the correctness set. The same
freezing therefore transfers the other uniform deficit statements of~\cite[Section~4]{SourceCovers2026}, its
Corollaries~4.4 and~4.7, with their hypotheses imposed on the geometry of~$\tau_0$.
\item Like the entropy hypothesis of Corollary~F$''$, hypothesis~(ii) concerns the full trace. Padding the input head's
motion with data-dependent detours that leave the computation unchanged can destroy it~\cite[Section~6.2]{BoundedNesting2026}.
\end{enumerate}
\end{remark}

\subsection{Block scale: an oblivious external-memory image}\label{sec:reductions-EM}

An external-memory simulation sees a machine's motion only when a head crosses a block boundary, and each crossing carries
a constant number of bits. Fixing these crossings therefore costs little input entropy, and it makes the simulation
oblivious. We use the following form of the external-memory model of Aggarwal and Vitter~\cite{AggarwalVitter1988}.

\begin{definition}[External-memory algorithms]\label{def:EM}
Fix a finite alphabet $\Sigma$ containing $0$, $1$ and a blank~$\sqcup$, and integers $B,m\ge1$ and $r\ge0$. A
\emph{block} is a word in $\Sigma^B$. A \emph{disk} is a map $D\colon\Z\to\Sigma^B$ whose values are blank blocks at all
but finitely many \emph{addresses}. The \emph{internal memory} consists of $m$ \emph{slots}, each holding one block, and a
\emph{register} of $r$ bits; its contents form the set $\mathcal M=(\Sigma^B)^m\times\{0,1\}^r$, and its size is
$mB\log_2\card\Sigma+r$ bits. An \emph{I/O} is a triple $(\sigma,i,a)$ with a direction
$\sigma\in\{\mathrm{read},\mathrm{write}\}$, a slot $i\in\{1,\dots,m\}$ and an address $a\in\Z$. An
\emph{external-memory algorithm} with these parameters is a map
\[
 \mathcal A\colon\mathcal M\to\{\mathrm{halt}\}\cup\bigl(\{\mathrm{read},\mathrm{write}\}\times\{1,\dots,m\}\times\Z
 \times\mathcal M\bigr).
\]
A \emph{run} from a disk~$D$ starts with blank slots and a zero register and repeats the following. If
$\mathcal A(\mu)=\mathrm{halt}$ for the current contents~$\mu$, the run stops. Otherwise
$\mathcal A(\mu)=(\sigma,i,a,\mu')$; a read replaces $\mu$ by $\mu'$ with slot~$i$ replaced by $D(a)$, and a write sets
$D(a)$ to slot~$i$ of $\mu'$ and replaces $\mu$ by~$\mu'$. The arbitrary map $\mu\mapsto\mu'$ models free internal
computation. The \emph{I/O sequence} of a run is the sequence of the I/Os $(\sigma,i,a)$ it performs, and their number is
its number of I/Os.

For transposition at dimension~$d$ the block size is $B=d\,b(d)$, one matrix row, so that $L=dB$. On input
$X\in\{0,1\}^L$ the initial disk holds row~$i$ of~$X$, that is the bits at input positions $iB,\dots,iB+B-1$, at
address~$i$ for $0\le i<d$, and is blank elsewhere. The algorithm is \emph{correct} on~$X$ if its run halts with row~$j$
of~$X^\T$ at address $d+j$ for every $0\le j<d$. It is \emph{oblivious on a set}~$\mathcal C$ of inputs if its runs on
all inputs in~$\mathcal C$ have the same I/O sequence.
\end{definition}

\begin{theorem}[Theorem EM: an oblivious external-memory image]\label{thm:EM}
\notformalized
Let $M$ be a machine with $k$ work tapes, state set~$Q$ and work alphabet~$\Gamma$, let $c>0$, and put
\[
 c_1=4(k+1)c+2,\qquad c_2=\bigl(2(k+1)c+1\bigr)\log_2(2k+4).
\]
There is $\gamma>0$, depending only on $k$, $c$ and $\card Q$, with the following property. Let $d\ge1$ be a dimension at
which $M$ transposes within $c(L+1)$ steps, and put $B=d\,b(d)$. Then there is an external-memory algorithm over the
alphabet~$\Sigma_M$ formed by $0$, $1$, the blank, the end marker and~$\Gamma$, with block size~$B$, $2k+3$ slots and a
register of at most $\gamma\log_2(2L)$ bits, such that:
\begin{enumerate}[label=(\roman*)]
\item it is correct on every input $X\in\{0,1\}^L$ and performs at most $c_1d$ I/Os on each;
\item it is oblivious on a set $\mathcal C$ of at least $2^{L-c_2d}$ inputs.
\end{enumerate}
\end{theorem}

\begin{informalproof}
Fix~$d$. Call the input tape tape~$0$ and the work tapes tapes $1,\dots,k$. Block~$j\in\Z$ of a tape consists of its cells
$jB,\dots,jB+B-1$.

\emph{Layout.} Block~$i$ of the input tape, for $0\le i<d$, is row~$i$ of the input, stored at address~$i$. The other
blocks of the input tape consist of end markers and are never stored. Block~$j$ of work tape~$h$ is stored at address
$2d+k\,z(j)+h-1$, where $z(j)=2j$ for $j\ge0$ and $z(j)=-2j-1$ for $j<0$; these addresses are distinct and at least
$2d$, and they are blank initially, as the work tapes are. The $n$-th group of $B$ consecutive output bits, $0\le n<d$, is
written to address $d+n$.

\emph{Windows.} Slots $2\tau+1$ and $2\tau+2$ hold the \emph{window} of tape~$\tau\in\{0,\dots,k\}$: its blocks $j_\tau$
and $j_\tau+1$, the lower one in slot $2\tau+1$. Slot $2k+3$ is an output buffer. The register holds the state of~$M$, the
head positions $x_\tau$ and window indices $j_\tau$, the number of output bits appended so far, and a flag for a pending
second I/O. We maintain three invariants: each window contains the current contents of its blocks; every work-tape block
outside its window has its current contents at its address; and the head of tape~$\tau$ lies in its window
$[j_\tau B,(j_\tau+2)B)$. Initially $j_\tau=-1$ for every~$\tau$. The work windows are then blank and the lower input slot
consists of end markers, both without an I/O, and the algorithm reads address~$0$ into slot~$2$; this first I/O is the
same for all inputs.

\emph{Simulation.} With the windows in memory, the algorithm simulates the steps of~$M$ without I/Os: the symbols under
the heads are in the windows, every work symbol written goes into its window, and every output bit appended goes into the
buffer. It stops simulating when, after a step, a head has left its window, or the buffer holds $B$ bits, or $M$ has
halted. These are the \emph{events}; the events of one step are processed in the order of the tapes, then the buffer.
\begin{itemize}
\item If the head of tape~$\tau$ has reached $(j_\tau+2)B$ (a \emph{right shift}), the algorithm writes the lower slot to the
address of block~$j_\tau$ if $\tau\ge1$, copies the upper slot into the lower one, fills the upper slot with block
$j_\tau+2$, and increases $j_\tau$ by one. The block is read from its address if $\tau\ge1$ or if $\tau=0$ and
$0\le j_\tau+2<d$; otherwise it consists of end markers and is produced without an I/O. A \emph{left shift}, when the head
has reached $j_\tau B-1$, is symmetric: write back block $j_\tau+1$ if $\tau\ge1$, move the lower slot up, fill the lower
slot with block $j_\tau-1$, and decrease~$j_\tau$.
\item If the buffer holds $B$ bits (a \emph{flush}), it is written to address $d+n$, where $n$ is the number of earlier
flushes, and emptied.
\item If $M$ has halted, the algorithm halts.
\end{itemize}
A shift performs at most two I/Os, and a flush one. The pending flag lets the second I/O of a shift be issued by the next
application of the map of \cref{def:EM}.

\emph{Correctness.} The invariants are preserved: a shift writes back the block that leaves the window and fills the
block that enters it with its current contents, and the head, which moves by at most one cell per step, is again inside
the window. Hence the simulation follows the run of~$M$ exactly. By transposition, $M$ appends exactly $L=dB$ bits, so
there are exactly $d$ flushes, all before~$M$ halts, and address $d+n$ receives output bits $nB,\dots,nB+B-1$, which form
row~$n$ of~$X^\T$. The algorithm is correct on every input.

\emph{Counting I/Os.} After a shift of tape~$\tau$ the head lies at offset $B$ (right shift) or $B-1$ (left shift) in the
new window, and at the start it lies at offset~$B$. In each case it is at distance at least~$B$ from both cells just
outside the window, so the next shift of that tape comes after at least $B$ moves of its head. A head moves at most $T$
times, where $T\le c(L+1)$ is the running time, so tape~$\tau$ shifts at most
\[
 \frac{c(L+1)}B=c\Bigl(d+\frac1B\Bigr)\le2cd
\]
times. The number of I/Os is therefore at most $1+d+2(k+1)\cdot2cd\le c_1d$, since $d\ge1$.

\emph{Obliviousness.} Record each event by a letter: $(\tau,+)$ or $(\tau,-)$ for a right or left shift of tape~$\tau$, and
$\phi$ for a flush. The \emph{event word} of a run is the sequence of these letters, a word over an alphabet of $2k+3$
letters, of length at most $(k+1)\cdot2cd+d=(2(k+1)c+1)d=:n_0$. The I/Os performed for an event are determined by its
letter and by the current window indices and flush count, which are determined by the earlier letters. Hence the I/O
sequence of a run is a function of its event word. There are at most $\sum_{n=0}^{\lfloor n_0\rfloor}(2k+3)^n\le(2k+4)^{n_0}$
words of length at most~$n_0$, so some event word is shared by a set $\mathcal C$ of at least
$2^L(2k+4)^{-n_0}=2^{L-c_2d}$ inputs, and the algorithm is oblivious on~$\mathcal C$.

\emph{Register.} The register holds the state of~$M$, a flag with at most $2k+1$ values, and $2k+3$ integers: the $x_\tau$
and $j_\tau$ and the output count. Each has absolute value at most $c(L+1)+2\le(2c+3)L$, since $\abs{x_\tau}\le T$ and
$\abs{j_\tau}\le\abs{x_\tau}+2$. As $\log_2(2L)\ge1$, the register needs at most $\gamma\log_2(2L)$ bits with
$\gamma=(2k+3)\bigl(\log_2(4c+7)+1\bigr)+\lceil\log_2\card Q\rceil+\lceil\log_2(2k+1)\rceil$.
\end{informalproof}

\Cref{thm:EM} thus places every linear-time transposer, adaptive or not, inside one concrete external-memory question:
a transposition with $O(d)$ I/Os and $O(B)$ bits of internal memory that is oblivious on a set of $2^{L-O(d)}$ inputs.
The constants $c_1,c_2$ do not depend on the alphabet, which affects only the memory size in bits.

\begin{corollary}[An external-memory route to the negation of (T)]\label{cor:EM}
\notformalized
Suppose that for every finite alphabet $\Sigma$ containing $0$, $1$ and a blank, and all constants $m\ge1$, $\gamma\ge0$
and $A\ge0$, there is $d_0$ such that for every $d\ge d_0$ the following holds at $L=d^2b(d)$ and $B=d\,b(d)$: every
external-memory algorithm over~$\Sigma$ with block size~$B$, $m$ slots and a register of at most $\gamma\log_2(2L)$ bits
that is correct and oblivious on some set of at least $2^{L-Ad}$ inputs performs more than $Ad$ I/Os on the inputs of that
set. Then (T) fails: no machine, with any number of work tapes, transposes in linear time.
\end{corollary}

\begin{informalproof}
Suppose that (T) holds, witnessed by a machine~$M$ with $k$ work tapes and a constant~$c$. Let $\gamma$, $c_1$ and $c_2$
be as in \cref{thm:EM}, and put $\Sigma=\Sigma_M$, $m=2k+3$ and $A=\max(c_1,c_2)$. Let $d_0$ be given by the hypothesis
for these constants, and let $d\ge d_0$. By \cref{thm:EM} there is an algorithm over~$\Sigma$ with block size $d\,b(d)$,
$m$ slots and a register of at most $\gamma\log_2(2L)$ bits that is correct on every input, oblivious on a set of at least
$2^{L-c_2d}\ge2^{L-Ad}$ inputs, and performs at most $c_1d\le Ad$ I/Os on each input. This contradicts the hypothesis.
\end{informalproof}

At block scale, then, deciding (T) negatively requires no hypothesis on adaptivity, nesting or coding beyond one
indivisibility-free I/O lower bound.

\begin{remark}[What the hypothesis of \cref{cor:EM} asks]\label{rem:EM-status}
The lower bounds of Aggarwal and Vitter for permuting and for matrix transposition~\cite{AggarwalVitter1988} assume that
records are indivisible: they are moved and copied but never combined. Under that assumption, transposition with a
constant number of blocks of internal memory needs more I/Os than linearly many, by a logarithmic factor; so a linear-time
machine yields, through \cref{thm:EM}, an algorithm that combines records, that is, a code. The hypothesis of
\cref{cor:EM} asks for a superlinear bound, of any order, without indivisibility, for algorithms that need only be correct and
oblivious on a set of $2^{L-O(d)}$ inputs, which need not be a product set. Farhadi, Hajiaghayi, Larsen and Shi prove
external-memory lower bounds for sorting without indivisibility, by network coding~\cite{FHLS2019}. They reduce
non-oblivious algorithms to a fixed I/O graph by a pigeonhole over I/O graphs, the step that produces~$\mathcal C$ in
\cref{thm:EM}. Their bounds are conditional on the conjecture that network coding gives no advantage over multicommodity
flow in undirected graphs, and that conjecture is false in general~\cite{BravermanHe2026}. We know of no unconditional
bound of the required kind. An indivisibility-free bound for this oblivious, large-set form of transposition is therefore
exactly what an impossibility proof needs at block scale. Adaptivity does not help there: fixing the I/O sequence costs
only $O(d)=o(L)$ bits of input entropy, so a machine's adaptivity can shape only its motion within blocks. There is no
converse to \cref{thm:EM}: internal computation is free and blocks are accessed at random, so an external-memory
transposition with $O(d)$ I/Os would not by itself give a linear-time machine.
\end{remark}
